\textbf{Proof.}
Work with normalized representatives of the projective classes in \(\mathcal B_\kappa\).  By prop:kappa-exhaustion and ass:kappa-feasibility, the quotient is a nonempty compact metrizable space.  Its block columns have squared norm \(n\), and \(\kappa\)-regularity gives
\[
 \lVert B\rVert_{\mathrm{op}}=\sqrt n,
 \qquad
 \lVert S_B^{-1}\rVert_{\mathrm{op}}\le(n\kappa)^{-1}.
 \tag{1}
\]
Exposure positivity makes \(\Pi^{-1}\) well defined; \(\Pi\) and \(\Omega\) are fixed.

Write \(q_a=F_a^{-1}\) and \(\widehat q_a=\widehat F_a^{-1}\).  If \(X_1,\ldots,X_m\) are independent, centered, and bounded in absolute value by one, independence gives
\[
 \mathbb E\!\left(\sum_{j=1}^mX_j\right)^4
 =\sum_j\mathbb EX_j^4+6\sum_{i<j}\mathbb EX_i^2\mathbb EX_j^2
 \le3m^2.
 \tag{2}
\]
Markov's inequality and Borel--Cantelli prove the bounded strong law.  Apply it in both pilot arms to outcomes, squared outcomes, and the indicators \(\mathbf1\{Y^P_{j,a}\le x\}\) simultaneously over rational \(x\).  Monotonicity gives empirical-distribution convergence at every continuity point.  The generalized-inverse squeeze, countability of the discontinuities of a monotone function, and dominated convergence on \([0,1]\) yield
\[
 \lVert\widehat q_a-q_a\rVert_{L^1(0,1)}\longrightarrow0,
 \qquad
 \widehat\mu_a\longrightarrow\mu_a,
 \qquad
 \widehat\sigma_a^2\longrightarrow\sigma_a^2
 \quad\text{almost surely}.
 \tag{3}
\]
Ass:unpaired-pilot supplies the two i.i.d. sequences, and the intersection of the two probability-one events still has probability one.

Since all quantiles lie in \([0,1]\),
\[
 \left|\int_0^1\widehat q_1(u)\widehat q_0(u)\,du
 -\int_0^1q_1(u)q_0(u)\,du\right|
 \le\lVert\widehat q_1-q_1\rVert_1+\lVert\widehat q_0-q_0\rVert_1,
 \tag{4}
\]
and the same inequality holds after replacing the second argument by \(1-u\).  Hence
\[
 \widehat c_L\longrightarrow c_L,
 \qquad \widehat c_U\longrightarrow c_U
 \quad\text{almost surely}.
 \tag{5}
\]
Ass:frechet-covariance-domain identifies these as the sharp endpoints used by the population criterion.  If \(Q_e=Q(c_e)\) and \(\widehat Q_e\) is its empirical analogue, def:moment-matrix and (3)--(5) imply
\[
 \rho_m:=\max_{e\in\{L,U\}}
 \lVert\widehat Q_e-Q_e\rVert_{\mathrm{op}}\longrightarrow0
 \quad\text{almost surely}.
 \tag{6}
\]

Put \(C_\nu=\lVert\Omega\Pi^{-1}\rVert_{\mathrm{op}}\) and \(A_{\mathrm{HT}}=\Pi^{-1}\Omega\Pi^{-1}\).  The inequality \(\lvert\operatorname{tr}M\rvert\le d\lVert M\rVert_{\mathrm{op}}\), def:gain-loss, and (1) give, uniformly in \(B\),
\[
 \max_e\lvert\widehat g_B(\widehat c_e)-g_B(c_e)\rvert
 \le\frac{2C_\nu^2}{n^2\kappa}\rho_m,
 \tag{7}
\]
whereas the \(2n\)-dimensional trace bound gives
\[
 \max_e\lvert\widehat h(\widehat c_e)-h(c_e)\rvert
 \le\frac2n\lVert A_{\mathrm{HT}}\rVert_{\mathrm{op}}\rho_m.
 \tag{8}
\]
Subtracting gain from uncontrolled variance and using the one-Lipschitz property of the maximum yields
\[
 \delta_m:=\sup_{B\in\mathcal B_\kappa}
 \lvert\widehat w(B)-w(B)\rvert\longrightarrow0
 \quad\text{almost surely}.
 \tag{9}
\]

All matrix operations are continuous on \(\mathcal B_\kappa\), because (1) keeps inversion away from singularity.  Thus \(w\) and \(\widehat w\) attain their minima, and
\[
 \lvert\widehat v_{\mathrm{mm}}-v_{\mathrm{mm}}\rvert\le\delta_m\longrightarrow0.
 \tag{10}
\]
Finally, take \(B_m\in\widehat{\mathcal A}_{\mathrm{mm}}\) and a convergent subsequence \(B_{m_j}\to B_\infty\).  For any \(B^\star\in\mathcal A_{\mathrm{mm}}\), empirical optimality gives
\[
 w(B_{m_j})\le\widehat w(B_{m_j})+\delta_{m_j}
 \le\widehat w(B^\star)+\delta_{m_j}
 \le v_{\mathrm{mm}}+2\delta_{m_j}.
 \tag{11}
\]
Continuity and \(w\ge v_{\mathrm{mm}}\) imply \(w(B_\infty)=v_{\mathrm{mm}}\).  Hence \(B_\infty\in\mathcal A_{\mathrm{mm}}\), proving the asserted outer limit inclusion. \(\square\)